\section{Fabre d’Olivet and the Myth of Blood}

\begin{quotex}
In the face of the chaos of modernity, the only salvation is form.
\flright{\textsc{Julius Evola}, motto to \emph{Il mito del sangue}}
\end{quotex}


\textbf{Julius Evola} regarded his books \textit{Il mito del sangue} (``The Myth of Blood'') and \textit{Sintesi di dottrina della razza} as two parts of a single work. The first volume is a systematic and critical study of various thinkers, most of them unknown or little read today, who maintained a biological basis for racial doctrine. His ultimate target was the National Socialists whose doctrines, he considered to be materialistic, debauched, and vulgar. In the second volume he proposed his own spiritual and transcendent view.

\textbf{Antoine Fabre d'Olivet} was a French Hermetist of the early 19\textsuperscript{th} century, who wrote important books on Pythagoras, music theory, the interpretation of the book of Genesis, and his most influential Philosophic History of the Human Race\footnote{\url{https://gornahoor.net/library/Hermeneutic\_interpretation\_of\_the\_origin.pdf}}. \textbf{Rene Guenon} and \textbf{Valentin Tomberg} have dealt with some of Fabre's ideas extensively. As for Evola, however, I have found just three references to him, all of them in Il mito del sangue. In the chapter on Gobineau, Evola mentions a writer named Virey who penned the book Natural History of the Human Race based on biological and genetic theories of race. Evola contrasts that with Fabre d'Olivet:

\begin{quotex}
Around the same time, here and there, various attempts were made to sketch our a spiritual classification. Thus, in 1824, with an obvious polemic reference to the work of Virey, Fabre d'Olivet's book of more than little interest, was published. \emph{The Philosophical History of the Human Race} contained a general racial classification scheme besides the attempt to identify the influence that, in a history starting with primordial times, each race would have exercised in its turn. The classifications are: the red or austral race, the yellow race, the black or sudeen race, the white or Boreal race. The most important aspect is that Fabre d'Olivet was the first, at that time, to hold to the theory of the remote Boreal or Hyperborean origin of the white race in the Arctic. For him, such a thesis has less of the character of a scientific hypothesis, than of the exposition of a traditional teaching, still conserved in certain very closed circles with which he was in contact.
\end{quotex}


Other writers then picked up on the theory of the Hyperborean origin, and Evola mentions, in particular, Teodoro Poesche, Karl Penka, and Ludwig Wilser. Unfortunately, we cannot get into their interesting speculations at this time. Suffice it to say that their ideas are ``polar'' opposite to the more contemporary generally accepted ``out of Africa'' scientific theory.

The final reference occurs in the chapter entitled The Arctic Myth, whose primary focus is \textbf{Hermann Wirth}. Evola contrasts the esoteric understanding of a Fabre d'Olivet to the attempts at a more ``scientific'' understanding. In particular, he is critical of those efforts that distort the esoteric understanding for political ends.

\begin{quotex}
As for Fabre d'Olivet, we have already said the ``arctic'', for him, is more than one of the many hypotheses of modern researchers. It corresponds instead to a knowledge of a ``traditional'' order, that has still been conserved among certain esoteric circles. Therefore, it has value independent of the efforts that those, such as Wirth, who have had an obscure intuition, and have tried to justify it with modern ``scientific'' methods and, above all, it has vale independent of the attempts of some racialists and even Wirth to utilize \emph{ad usum delphini}, that is, for more or less contingent political goals.
\end{quotex}



\flrightit{Posted on 2011-08-03 by Cologero}

\begin{center}* * *\end{center}

\begin{footnotesize}
\begin{sffamily}
\texttt{Doctor Bones on 2011-08-04 at 06:21 said:}

``The classifications are: the red or southern race, the yellow race, the black race, the white or Boreal race.''

This is an ancient scheme of the human races, corresponding to the four cardinal points. I knew it to be thus: North/White, South/Black, East/Yellow, West/Red. Notice the correspondence of this scheme to actual racial distribution amongst the continents. Any comments on why south is connected to the red race in the quote above?

\hfill

\texttt{Cologero on 2011-08-04 at 12:19 said:}

Good point. I updated the post to use the terms untranslated (e.g., ``austral'' instead of ``southern''). This is the same practice used in the English translation of the French original. I put Fabre d'Olivet's book in the library so you can compare how Fabre d'Olivet himself uses these terms.

\hfill

\texttt{HOO on 2011-08-07 at 09:06 said:}

Anyone with a mind retaining enough clarity will see Evola is not opposed to the Political use of the Arctic or Nordic Myth, quite the contrary, but is merely opposed to its misuse by profane politics. Of course the Myth of Origins has been very Political to all superior tribes and civilizations, as is seen how they trace their bloodlines.

Evola himself concludes his chapter on Wirth from MS with these words:

... the period preceding liberalism and scientism was characterized by three fundamental ideas:

1. The equality of the human race

2. Nordic barbarianism and the origin of every civilisation from the East

3. The Hebraic origin of monotheism

Wirth struck down or overthrew those ideas with these three: ... :

1. Humanity is differentiated into distinct races

2. Civilisation did not come from the East, but from the North

3. It was the Nordics who would have known a higher monotheistic religion well before the Hebrews

\hfill

\texttt{Cologero on 2011-08-07 at 10:43 said:}

Precisely, Hoo, but ``use'' is not the correct word. Political Power flows out of Spiritual Authority, so it is not really the ``use'' of a myth for political ends. That may be why Evola considers the emphasis on so-called scientific evidence to be misguided; once the spiritual roots are known, then the suitable political arrangement follows. It is no longer a case of ``we \emph{should} organize in such and such a way, because our genetic heritage is such and such.

I'm rushing to get more translated, for example, the entire chapter on Wirth. It seems Evola is speaking from beyond the grave, issuing a warning about the repetition of past mistakes. A false biological view of race led to an equally false anti-racism. In the chapter under discussion, Evola points to the Nat Soc policy of requiring a certain ethnic heritage as a prerequisite to citizenship. So, today, as an over-reaction to those policies, we see that, throughout Europe, ethnic heritage is of no relevance to citizenship. So what should be the basis of citizenship. That is the reason we opened up the question about \textit{what is a nation}\footnote{\url{http://www.gornahoor.net/?p=2775}}. As Evola's critique of the lack of spirit of the Scandinavians of his time shows, this is a spiritual, not a biological, question.

We've pointed out how the Germanics took the husk of the Roman Empire and created great civilization based on a military aristocracy; they also created a spiritual system that unified Europe from Portugal to the Urals. We can also look at the accomplishments of the Normans, the sons of Norway, a millennium ago. What has changed? The current narrative is that the Northern Europeans have high IQs, so they are brilliant engineers and keep the trains running on time. Is that all there is?

\hfill

\texttt{Cologero on 2011-08-07 at 10:56 said:}

From Il mito del sangue:

\begin{quotex}\footnotesize

In any case, it is necessary to clearly distinguish the value and the meaning of the arctic, or Hyperborean, thesis in itself from Wirth's arbitrary personal adaptations, because the plane on which it belongs is quite different and has a totally other dignity than these reconstructions of contemporary researchers, reconstructions, nevertheless, not without interest as signs and as obscure presentiments oscura presentimento of a truth.
\end{quotex}


\hfill

\texttt{HOO on 2011-08-07 at 11:34 said:}

The Urnordisch Myth is multilayered as any myth. It tells a ´story´ or delivers a truth, being and power, on all levels; from the spiritual-worldview down to the psychological and historical levels (or even more contingent levels). For example, even in fantasy stories, there are splendiferous races, being obscure presentiments of the Primordial Race.

´We have referred to a primordial Nordic tradition. It is not a myth, it is our truth.´ Evola 

\hfill

\texttt{Cologero on 2011-08-07 at 17:07 said:}

Yes, Hoo, but we are at the point where we can no longer be satisfied with ``obscure presentiments'' nor the fantasy stories made known by the likes of Godwin or Goodrick-Clarke. Even less so the false path of a counter-Tradition. Quo vadis?

\hfill

\texttt{HOO on 2011-08-08 at 05:28 said:}

Of course no one who isn't a lower-grade pashu is satisfied with mere ``obscure presentiments'', nor fantasy stories, or mere scholarliness.

Where am I going? Preferably in the direction of a ´kind of `immanent transcendence': by the direct presence of non-being (in its positive sense of supra-ontological essentiality) within being, of the infinitely remote (the `Sky') in what is close, and of what is beyond nature within nature.´


\hfill

\end{sffamily}
\end{footnotesize}

